\chapter{Dirac 理论、Foldy--Wouthuysen 展开与原子精细结构}\label{ch:spin-orbit-thomas}

原子精细结构若从“电子静止系中的磁场”开始，会把一个相对论量子问题倒置成半经典猜测。正确主线是：先由 Lorentz 对称性建立 Dirac 母方程，再对正能量子空间作受控低能展开；Pauli 磁矩、相对论动能、Darwin 项和自旋轨道耦合会在同一 Foldy--Wouthuysen（FW）块对角化中同时出现。Thomas--Wigner 进动随后作为同一系数的经典几何解释，而不是决定系数的替代推导。

\section{Lorentz 表示、Clifford 代数与 Dirac 最小耦合}\label{sec:dirac-lorentz}

取 Minkowski 度规
\[
\eta_{\mu\nu}=\operatorname{diag}(1,-1,-1,-1),
\qquad
x^\mu=(ct,\mathbf r),
\qquad
A_\mu=(\Phi/c,-\mathbf A).
\]
Lorentz 代数生成元 $M^{\mu\nu}=-M^{\nu\mu}$ 满足
\begin{equation}
[M^{\mu\nu},M^{\rho\sigma}]
=\ii\hbar
\left(
\eta^{\nu\rho}M^{\mu\sigma}
-\eta^{\mu\rho}M^{\nu\sigma}
-\eta^{\nu\sigma}M^{\mu\rho}
+\eta^{\mu\sigma}M^{\nu\rho}
\right).
\label{eq:rel-lorentz-algebra}
\end{equation}
旋量表示由 Clifford 代数
\begin{equation}
\{\gamma^\mu,\gamma^\nu\}=2\eta^{\mu\nu}\mathbb I,
\qquad
\sigma^{\mu\nu}\equiv\frac{\ii}{2}[\gamma^\mu,\gamma^\nu]
\label{eq:clifford-algebra}
\end{equation}
实现，其 Lorentz 生成元为
\begin{equation}
S^{\mu\nu}=\frac{\hbar}{2}\sigma^{\mu\nu}
=\frac{\ii\hbar}{4}[\gamma^\mu,\gamma^\nu].
\label{eq:spinor-lorentz-generator}
\end{equation}

对质量 $m$、电荷 $Q$ 的自旋 $1/2$ 粒子，采用四维规范约定
\[
A_\mu' = A_\mu-\partial_\mu\chi,
\qquad
\psi'=\exp\!\left(\frac{\ii Q\chi}{\hbar}\right)\psi,
\]
并定义
\[
D_\mu=\partial_\mu+\frac{\ii Q}{\hbar}A_\mu.
\]
则 $D_\mu'\psi'=\exp(\ii Q\chi/\hbar)D_\mu\psi$，Dirac 方程写成
\begin{equation}
\boxed{
\left(\ii\hbar c\gamma^\mu D_\mu-mc^2\right)\psi=0.
}
\label{eq:dirac}
\end{equation}
Schr\"odinger 形式为
\begin{equation}
\ii\hbar\partial_t\psi
=H_D\psi,
\qquad
H_D=c\boldsymbol\alpha\cdot\boldsymbol\pi
+\beta mc^2+Q\Phi,
\qquad
\boldsymbol\pi=\mathbf p-Q\mathbf A,
\label{eq:dirac-hamiltonian-em}
\end{equation}
其中 $\alpha^i=\gamma^0\gamma^i$、$\beta=\gamma^0$。守恒流
\begin{equation}
j^\mu=c\bar\psi\gamma^\mu\psi,
\qquad
\partial_\mu j^\mu=0,
\qquad
\bar\psi=\psi^\dagger\gamma^0
\label{eq:dirac-current}
\end{equation}
给出正定密度 $j^0=c\psi^\dagger\psi$。

\begin{derivation}[从 Clifford 代数到 Lorentz 旋量生成元]
无穷小 Lorentz 变换写成
\[
\Lambda^\rho{}_{\sigma}
=\delta^\rho{}_{\sigma}+\omega^\rho{}_{\sigma}+O(\omega^2),
\qquad
\omega_{\mu\nu}=-\omega_{\nu\mu}.
\]
旋量变换取
\[
\mathcal S(\Lambda)
=\mathbb I-\frac{\ii}{2\hbar}\omega_{\mu\nu}S^{\mu\nu}+O(\omega^2).
\]
由 Clifford 代数直接计算
\[
[S^{\mu\nu},\gamma^\rho]
=\ii\hbar
\left(\eta^{\nu\rho}\gamma^\mu-\eta^{\mu\rho}\gamma^\nu\right).
\]
于是
\[
\begin{aligned}
\mathcal S^{-1}\gamma^\rho\mathcal S
&=\gamma^\rho
+\frac{\ii}{2\hbar}\omega_{\mu\nu}
[S^{\mu\nu},\gamma^\rho]+O(\omega^2)\\
&=\gamma^\rho+\omega^\rho{}_{\sigma}\gamma^\sigma+O(\omega^2),
\end{aligned}
\]
恰好得到四矢量变换律。再用 Jacobi 恒等式可得
\[
[S^{\mu\nu},S^{\rho\sigma}]
=\ii\hbar
\left(
\eta^{\nu\rho}S^{\mu\sigma}
-\eta^{\mu\rho}S^{\nu\sigma}
-\eta^{\nu\sigma}S^{\mu\rho}
+\eta^{\mu\sigma}S^{\nu\rho}
\right),
\]
与\eqnrefs{eq:rel-lorentz-algebra}同构。因此 $\gamma^\mu$ 与 $S^{\mu\nu}$ 是 Lorentz 旋量表示的结构，而不是仅为“线性化能量色散”临时引入的矩阵。

\textbf{物理意义。}Dirac $\gamma$ 矩阵不只是线性化 $\sqrt{p^2c^2+m^2c^4}$ 的工具——它们携带 Lorentz 群的旋量表示，这是电子自旋的群论起源。$S^{\mu\nu}=\frac14[\gamma^\mu,\gamma^\nu]$ 生成 Lorentz 变换在旋量空间的作用，自旋算符 $\mathbf S=\frac12\boldsymbol\Sigma$ 从中自然产生。实例：Weyl 旋量（手征表示）描述中微子，Dirac 旋量描述电子，Majorana 旋量描述中微子若其是自身反粒子。
\end{derivation}

\section{Foldy--Wouthuysen 块对角化与 Breit--Pauli 结构}

原子束缚态满足 $|E-m_ec^2|\ll m_ec^2$。需要的是 Dirac 正能量子空间上的系统低能展开，而不是直接“删除小分量”。把静态 Dirac Hamiltonian 写成
\[
H_D=\beta mc^2+\mathcal E+\mathcal O,
\qquad
[\beta,\mathcal E]=0,
\qquad
\{\beta,\mathcal O\}=0,
\]
其中 $\mathcal E=Q\Phi$、$\mathcal O=c\boldsymbol\alpha\cdot\boldsymbol\pi$。对原子精细结构采用静电主导计数
\[
\frac{p}{mc}=O(v/c)\ll1,
\qquad
\frac{|Q\Phi|}{mc^2}=O(v^2/c^2),
\]
并只保留外磁场的领先 Pauli 耦合。FW 变换给出正能量块
\begin{equation}
\boxed{
\begin{aligned}
H_{\rm FW}^{(+)}
={}&mc^2+Q\Phi
+\frac{\boldsymbol\pi^2}{2m}
-\frac{\boldsymbol\pi^4}{8m^3c^2}
-\frac{Q\hbar}{2m}\boldsymbol\sigma\cdot\mathbf B\\
&-\frac{Q\hbar^2}{8m^2c^2}\nabla\cdot\mathbf E
-\frac{Q\hbar}{8m^2c^2}
\boldsymbol\sigma\cdot
\left(\mathbf E\times\boldsymbol\pi-\boldsymbol\pi\times\mathbf E\right)
+O(c^{-4}).
\end{aligned}
}
\label{eq:fw-general-order-c2}
\end{equation}
这些项分别是非相对论动能、相对论动能修正、Pauli 磁矩、Darwin 项和自旋轨道项。对任意强静磁场，$\mathcal O^4$ 中还存在同阶相对论磁修正；\eqnrefs{eq:fw-general-order-c2}针对的是原子 Coulomb 场主导的精细结构计数。

\begin{derivation}[FW even--odd 消元到相对阶 \(1/c^2\)]
第一步取生成元
\[
S_1=-\frac{\ii}{2mc^2}\beta\mathcal O.
\]
Baker--Campbell--Hausdorff 展开
\[
H'=
\ee^{\ii S_1}H_D\ee^{-\ii S_1}
=H_D+\ii[S_1,H_D]-\frac12[S_1,[S_1,H_D]]+\cdots
\]
中，$\ii[S_1,\beta mc^2]$ 恰好消去最低阶 odd 块。重复消元并保留目标阶，even 部分为
\begin{equation}
H_{\rm FW}
=\beta mc^2+\mathcal E
+\frac{\beta\mathcal O^2}{2mc^2}
-\frac{\beta\mathcal O^4}{8m^3c^6}
-\frac{[\mathcal O,[\mathcal O,\mathcal E]]}{8m^2c^4}
+O(c^{-4}).
\label{eq:fw-even-odd-master}
\end{equation}
利用
\[
[\pi_i,\pi_j]=\ii Q\hbar\epsilon_{ijk}B_k,
\qquad
\alpha_i\alpha_j=\delta_{ij}+\ii\epsilon_{ijk}\Sigma_k,
\]
得到
\[
\mathcal O^2
=c^2(\boldsymbol\pi^2-Q\hbar\boldsymbol\Sigma\cdot\mathbf B).
\]
对静电标势，$[\pi_i,Q\Phi]=\ii\hbar Q E_i$，进一步计算
\[
[\mathcal O,[\mathcal O,Q\Phi]]
=Q\hbar^2c^2\nabla\cdot\mathbf E
+Q\hbar c^2\boldsymbol\Sigma\cdot
(\mathbf E\times\boldsymbol\pi-\boldsymbol\pi\times\mathbf E).
\]
在正能量块取 $\boldsymbol\Sigma\to\boldsymbol\sigma$。静电主导计数下，$\mathcal O^4$ 在本阶只需保留 $c^4\boldsymbol\pi^4$；若磁场不属于弱修正，则 $\{\pi^2,\boldsymbol\sigma\cdot\mathbf B\}$ 等项也必须同时保留。代回\eqnrefs{eq:fw-even-odd-master}即得\eqnrefs{eq:fw-general-order-c2}。

\textbf{物理意义。}Foldy--Wouthuysen 变换把 Dirac 方程分解为正能量（电子）和负能量（正电子）两个块对角部分，消除了奇算符（混合正负能量的项）。$1/c^2$ 阶给出 Pauli 方程加 Darwin 项（接触势）和自旋-轨道耦合 $\frac{e\hbar}{4m^2c^2}\boldsymbol\sigma\cdot(\mathbf E\times\mathbf p)$——后者解释原子光谱的精细结构。实例：氢原子 $2p_{3/2}-2p_{1/2}$ 分裂 $\Delta E=\frac{(Z\alpha)^4mc^2}{32}$ 正来自此自旋-轨道项。
\end{derivation}

对纯静电中心势 $U(r)=Q\Phi(r)$，$\mathbf B=0$ 且 $\mathbf E=-\nabla\Phi$。于是
\begin{equation}
H_{\rm so}
=\frac{1}{2m^2c^2r}\frac{\dd U}{\dd r}\,\mathbf L\cdot\mathbf S,
\qquad
H_{\rm Darwin}
=\frac{\hbar^2}{8m^2c^2}\nabla^2U.
\label{eq:pauli-from-dirac}
\end{equation}

\begin{derivation}[从 FW 自旋-轨道项到中心势形式]
\emph{思路：}FW 展开给出的自旋-轨道耦合 \(H_{\rm so}=\frac{e\hbar}{4m^2c^2}\boldsymbol\sigma\cdot(\mathbf E\times\mathbf p)\) 中电场 \(\mathbf E\) 是一般矢量。对中心势 \(U(r)\)，电场沿径向，叉乘 \(\mathbf E\times\mathbf p\) 自然给出轨道角动量 \(\mathbf L=\mathbf r\times\mathbf p\)。Darwin 项来自 FW 展开中 \(\nabla\cdot\mathbf E\) 的接触项。

\emph{第一步：中心势的电场。}对中心势 \(U(r)=Q\Phi(r)\)，电场为 \(\mathbf E=-\nabla\Phi=-\frac{1}{Q}\frac{dU}{dr}\hat{\mathbf r}\)。这里 \(\hat{\mathbf r}=\mathbf r/r\) 是径向单位矢量。对 Coulomb 势 \(U=-Ze^2/(4\pi\epsilon_0 r)\)，\(\frac{dU}{dr}=Ze^2/(4\pi\epsilon_0 r^2)\)，电场 \(\mathbf E=\frac{Ze}{4\pi\epsilon_0 r^2}\hat{\mathbf r}\) 指向远离核的方向（对正核 \(Z>0\)）。

\emph{第二步：叉乘化为角动量。}关键代数步骤：
\begin{equation*}
\mathbf E\times\mathbf p=-\frac{1}{Qr}\frac{dU}{dr}\,\mathbf r\times\mathbf p=-\frac{1}{Qr}\frac{dU}{dr}\,\mathbf L,
\end{equation*}
其中 \(\mathbf L=\mathbf r\times\mathbf p\) 是轨道角动量。这一步把"电场叉乘动量"转化为"径向力叉乘动量"=角动量乘径向系数——物理上，自旋-轨道耦合是自旋与轨道运动的相互作用，中心势下轨道运动由 \(\mathbf L\) 表征。

\emph{第三步：代入并化简。}代入 \(H_{\rm so}\) 并利用 \(\mathbf S=\frac{\hbar}{2}\boldsymbol\sigma\)：
\begin{equation*}
H_{\rm so}=\frac{e\hbar}{4m^2c^2}\boldsymbol\sigma\cdot\left(-\frac{1}{Qr}\frac{dU}{dr}\mathbf L\right)
=-\frac{e}{2m^2c^2Qr}\frac{dU}{dr}\,\mathbf S\cdot\mathbf L.
\end{equation*}
对电子 \(Q=-e\)，负号消去，得
\begin{equation*}
H_{\rm so}=\frac{1}{2m^2c^2r}\frac{dU}{dr}\,\mathbf L\cdot\mathbf S,
\end{equation*}
即\eqnrefs{eq:pauli-from-dirac}的第一式。系数中的 \(\frac{1}{2}\) 就是 Thomas 因子——由 Dirac/FW 结构自动给出，无需手工插入。历史上，非相对论推导先得到系数 1（遗漏 Thomas 进动），再手工减半修正；Dirac 方程自动给出正确的 \(\frac{1}{2}\)，是其正确性的一个标志。

\emph{第四步：Darwin 项。}Darwin 项 \(H_{\rm Darwin}=\frac{\hbar^2}{8m^2c^2}\nabla^2 U\) 来自 FW 展开中 \(\nabla\cdot\mathbf E\) 的接触项。对 Coulomb 势 \(U=-Ze^2/(4\pi\epsilon_0 r)\)，利用 \(\nabla^2(1/r)=-4\pi\delta^{(3)}(\mathbf r)\)：
\begin{equation*}
\nabla^2 U=\frac{Ze^2}{\epsilon_0}\delta^{(3)}(\mathbf r),
\end{equation*}
故 \(H_{\rm Darwin}\propto\delta^{(3)}(\mathbf r)\)，只对 \(s\) 波（\(L=0\)，\(|\psi(0)|^2\ne0\)）有贡献。物理上，Darwin 项反映电子不是点粒子而是有 \(10^{-13}\,\mathrm{cm}\) 尺度的 Zitterbewegung（颤动），在核附近"抹平"了 Coulomb 奇点。

\emph{第五步：氢原子精细结构。}对氢原子 \(U=-Ze^2/(4\pi\epsilon_0 r)\)，\(dU/dr=Ze^2/(4\pi\epsilon_0 r^2)\)，故
\begin{equation*}
H_{\rm so}=\frac{Ze^2}{8\pi\epsilon_0 m^2c^2 r^3}\mathbf L\cdot\mathbf S.
\end{equation*}
由 \(\mathbf J=\mathbf L+\mathbf S\) 得 \(\mathbf L\cdot\mathbf S=\frac{1}{2}(J^2-L^2-S^2)\)，本征值为 \(\langle\mathbf L\cdot\mathbf S\rangle=\frac{\hbar^2}{2}[J(J+1)-L(L+1)-S(S+1)]\)。对 \(2p\) 态（\(L=1\)，\(S=\frac{1}{2}\)）：\(J=\frac{3}{2}\) 给出 \(\langle\mathbf L\cdot\mathbf S\rangle=\frac{\hbar^2}{2}\)，\(J=\frac{1}{2}\) 给出 \(\langle\mathbf L\cdot\mathbf S\rangle=-\hbar^2\)。\(2p_{3/2}\) 与 \(2p_{1/2}\) 的分裂 \(\Delta E=\frac{(Z\alpha)^4 mc^2}{32}\)，对氢（\(Z=1\)）约 \(4.5\times10^{-5}\,\mathrm{eV}\)，与实验精确吻合——这是 Dirac 方程最早期也是最优美的验证之一。精细结构分裂的一般公式为
\begin{equation*}
\Delta E_{n,j}=\frac{(Z\alpha)^4 mc^2}{2n^3}\left(\frac{1}{j+\frac12}-\frac{3}{4n}\right),
\end{equation*}
对 \(n=2\)，\(j=1/2\) 和 \(j=3/2\) 给出上述分裂。






\end{derivation}

对氢样电子 $Q=-e$、$U(r)=-Ze^2/(4\pi\epsilon_0r)$，
\[
H_{\rm so}
=\frac{Ze^2}{8\pi\epsilon_0m_e^2c^2r^3}\,\mathbf L\cdot\mathbf S.
\]
Darwin 项只对在原点非零的 $s$ 波径向分量有贡献。自旋轨道项中的 $1/2$ 已由 Dirac/FW 结构固定，无需再手工插入 Thomas 因子。

\section{LS、jj 与中间耦合：精细结构和 Zeeman 结构}

\chapref{ch:angular-momentum}已经建立一般角动量耦合。对给定 $L,S,J$，
\begin{equation}
\langle\mathbf L\cdot\mathbf S\rangle
=\frac{\hbar^2}{2}
[J(J+1)-L(L+1)-S(S+1)].
\label{eq:ls-expectation}
\end{equation}
若同一 $LS$ term 中精细结构可以写成
\[
H_{\rm fs}=A_{LS}\,\mathbf L\cdot\mathbf S/\hbar^2,
\]
则
\[
E_J
=E_{LS}+\frac{A_{LS}}{2}
[J(J+1)-L(L+1)-S(S+1)],
\]
从而
\begin{equation}
E_J-E_{J-1}=A_{LS}J.
\label{eq:lande-interval-energy}
\end{equation}
这就是 Land\'e interval rule。它要求 LS coupling 是良好近似；若自旋轨道与电子间静电分裂同阶，必须在固定 $J$ 子空间内整体对角化，即 intermediate coupling。

弱磁场 $\mathbf B=B\mathbf e_z$ 中，
\[
H_Z=\frac{\mu_B}{\hbar}(\mathbf L+g_s\mathbf S)\cdot\mathbf B,
\qquad
\mu_B=\frac{e\hbar}{2m_e}.
\]
在 $|LSJM\rangle$ 中作矢量投影，得到
\begin{equation}
\Delta E_Z=\mu_Bg_JMB,
\qquad
g_J
=1+\frac{J(J+1)+S(S+1)-L(L+1)}{2J(J+1)}(g_s-1).
\label{eq:lande-g-factor}
\end{equation}
取 Dirac 树级 $g_s=2$ 得常用 Land\'e 因子。若 Zeeman 能量超过自旋轨道耦合，$J$ 不再是好量子数，系统趋向 $m_L,m_S$ 近似独立的 Paschen--Back 极限。

\begin{derivation}[Land\'e 因子来自向总角动量的投影]
在固定 $J$ 不可约子空间内，由 Wigner--Eckart 定理，任意矢量算符的一阶对角矩阵元都必须平行于 $\mathbf J$，故可写
\begin{equation*}
\langle\mathbf L\rangle = a_L\,\langle\mathbf J\rangle,
\qquad
\langle\mathbf S\rangle = a_S\,\langle\mathbf J\rangle,
\end{equation*}
其中 $a_L,a_S$ 为待定标量。由 $\mathbf J=\mathbf L+\mathbf S$，于是
\begin{equation*}
\langle\mathbf J\rangle
= \langle\mathbf L\rangle+\langle\mathbf S\rangle
= (a_L+a_S)\,\langle\mathbf J\rangle,
\end{equation*}
因此 $a_L+a_S=1$。为求 $a_L$，将上式两边点乘 $\mathbf J$ 并取对角元，
\begin{equation*}
\langle\mathbf L\cdot\mathbf J\rangle
= a_L\,\langle J^2\rangle.
\end{equation*}
由 $\mathbf J=\mathbf L+\mathbf S$ 平方得
\begin{equation*}
\mathbf L\cdot\mathbf J
= \frac12\bigl(J^2+L^2-S^2\bigr),
\qquad
\mathbf S\cdot\mathbf J
= \frac12\bigl(J^2+S^2-L^2\bigr),
\end{equation*}
代回并取 $|LSJM\rangle$ 对角元，
\begin{equation*}
\frac12\bigl[J(J+1)+L(L+1)-S(S+1)\bigr]\hbar^2
= a_L\,J(J+1)\hbar^2,
\end{equation*}
于是
\begin{equation*}
a_L
= \frac{J(J+1)+L(L+1)-S(S+1)}{2J(J+1)}.
\end{equation*}
同理对 $\mathbf S$ 重复点乘 $\mathbf J$ 的步骤，
\begin{equation*}
\langle\mathbf S\cdot\mathbf J\rangle
= a_S\,\langle J^2\rangle,
\end{equation*}
代入第二个恒等式并取对角元，
\begin{equation*}
a_S
= \frac{J(J+1)+S(S+1)-L(L+1)}{2J(J+1)}.
\end{equation*}
因此
\begin{equation*}
\langle\mathbf L\rangle
= \frac{J(J+1)+L(L+1)-S(S+1)}{2J(J+1)}\,\langle\mathbf J\rangle,
\qquad
\langle\mathbf S\rangle
= \frac{J(J+1)+S(S+1)-L(L+1)}{2J(J+1)}\,\langle\mathbf J\rangle.
\end{equation*}
由磁矩 $\boldsymbol\mu=-(\mu_B/\hbar)(\mathbf L+g_s\mathbf S)$ 与 Zeeman 微扰 $H_Z=-\boldsymbol\mu\cdot\mathbf B$，代回投影式，
\begin{equation*}
\langle H_Z\rangle
= \frac{\mu_B}{\hbar}\,\bigl(a_L+g_s a_S\bigr)\,\langle J_z\rangle B.
\end{equation*}
取 $\langle J_z\rangle=\hbar M$，整理得
\begin{equation*}
\Delta E_Z = \mu_B g_J M B,
\qquad
g_J = a_L+g_s a_S
= 1+\frac{J(J+1)+S(S+1)-L(L+1)}{2J(J+1)}\,(g_s-1),
\end{equation*}
即\eqnrefs{eq:lande-g-factor}。

\textbf{物理意义。}Land\'e 因子 $g_J$ 把原子磁矩与总角动量 $J$ 联系起来，是实验观测 Zeeman 分裂 $\Delta E=\mu_B g_J M B$ 直接读出的量。不同 $LS$ 耦合态的 $g_J$ 值是光谱学中鉴定原子态的指纹：测量 Zeeman 分裂的条数和间距即可反推 $L,S,J$。

\textbf{典型实例。}碱金属基态 ${}^2S_{1/2}$：$L=0,S=1/2,J=1/2$，$g_J=2$（纯自旋）。钠 $D$ 线激发态 ${}^2P_{1/2}$：$L=1,S=1/2,J=1/2$，$g_J=2/3$。${}^2P_{3/2}$：$J=3/2$，$g_J=4/3$。${}^1D_2$（单态）：$S=0$，$g_J=1$（纯轨道）。这些值可直接由 $g_J=1+(g_s-1)[J(J+1)+S(S+1)-L(L+1)]/[2J(J+1)]$ 取 $g_s=2$ 验证。$g_J$ 偏离 2 的程度正是轨道角动量贡献的度量。
\end{derivation}

\section{Thomas--Wigner 进动与 Dirac 低能极限}

两个不共线 Lorentz boost 的乘积并不是纯 boost，而会附带空间旋转。对速度 $\mathbf v(t)$ 的粒子，相邻瞬时共动系之间的 Thomas--Wigner 角速度为
\begin{equation}
\boldsymbol\Omega_T
=\frac{\gamma^2}{\gamma+1}
\frac{\mathbf a\times\mathbf v}{c^2}.
\label{eq:thomas-exact}
\end{equation}
在 $v\ll c$ 时
\begin{equation}
\boldsymbol\Omega_T
=\frac{1}{2c^2}\mathbf a\times\mathbf v
+O(v^4/c^4).
\label{eq:thomas-nonrel}
\end{equation}
电子在中心电场中的瞬时静止系会看到运动诱导磁场；若只计算该磁场引起的 Larmor 进动，会得到两倍于\eqnrefs{eq:pauli-from-dirac}的自旋轨道系数。共动坐标系自身还按\eqnrefs{eq:thomas-nonrel}旋转，恰好把多出的因子减半。因此 Thomas 因子是 boost 非交换性的经典几何表现，与 Dirac/FW 的量子结果一致。

自由 Dirac Hamiltonian 还给出速度算符 $\dot{\mathbf x}=c\boldsymbol\alpha$。Heisenberg 方程可解为
\begin{equation}
\boldsymbol\alpha(t)
=c\mathbf pH_D^{-1}
+\left[\boldsymbol\alpha(0)-c\mathbf pH_D^{-1}\right]
\ee^{-2\ii H_Dt/\hbar},
\label{eq:alpha-zitter}
\end{equation}
从而位置包含频率约 $2mc^2/\hbar$ 的正负能量干涉项，即 Zitterbewegung。FW 表象把正负能量子空间在给定阶数上解耦后，这一快速干涉不再作为低能正能量波包的独立运动项。Pauli 理论因此应理解为 Dirac 谱子空间的受控块对角化，而不是“删去两个分量”。

\section{精细结构、兰姆移位与实验数值}
\label{sec:rel-fine-structure-experiment}

前几节从 Dirac 方程经 FW 变换导出了自旋轨道耦合、Darwin 项和相对论修正。本节把这些修正代入氢原子和碱金属，给出可与光谱实验直接对比的数值。

接下来考察氢原子精细结构。

Dirac 方程对氢原子的精确能级为
\[
 E_{n,j}=mc^2\left[1+\frac{(Z\alpha)^2}{\left(n-j-\frac12+\sqrt{(j+\frac12)^2-(Z\alpha)^2}\right)^2}\right]^{-1/2},
\]
展开到 $(Z\alpha)^4$ 阶（减去静能 $mc^2$）得精细结构公式
\[
 E_{n,j}^{(0)}+E_{\rm FS}
 =-\frac{R_\infty hc}{n^2}\left[1+\frac{(Z\alpha)^2}{n}\left(\frac1{j+\frac12}-\frac{3}{4n}\right)\right].
\]
对氢原子 $Z=1$、$\alpha\approx1/137$，典型精细结构分裂数值：
\begin{itemize}
 \item $2p_{3/2}-2p_{1/2}$：$\Delta E=4.53\times10^{-5}\,\mathrm{eV}$，对应频率 $10.97\,\mathrm{GHz}$，波长差 $\Delta\lambda\approx0.006\,\mathrm{nm}$。
 \item $2s_{1/2}-2p_{1/2}$：Dirac 理论给出 $\Delta E=0$（简并），但 QED 修正（兰姆移位）破缺此简并。
 \item $3p_{3/2}-3p_{1/2}$：$\Delta E=1.34\times10^{-5}\,\mathrm{eV}$，对应 $3.24\,\mathrm{GHz}$。
\end{itemize}
精细结构常数 $\alpha=e^2/(4\pi\varepsilon_0\hbar c)\approx1/137.036$ 是这些分裂的无量纲控制参数：$\Delta E_{\rm FS}/E_{\rm Bohr}\sim(Z\alpha)^2\sim5\times10^{-5}$。

接下来考察兰姆移位：QED 对 Dirac 理论的修正。

1947 年 Lamb 和 Retherford 用微波共振实验测得氢 $2s_{1/2}$ 比 $2p_{1/2}$ 高
\[
 \Delta E_{\rm Lamb}=E(2s_{1/2})-E(2p_{1/2})=1057.8\,\mathrm{MHz}\times h=4.37\,\mathrm{\mu eV}.
\]
这一分裂在 Dirac 理论中严格为零，其物理来源是真空涨落（电子与真空电磁场零点涨落耦合）和电子自能修正。Bethe 的初步估算给出 $\Delta E_{\rm Lamb}\approx\frac{4\alpha^5mc^2}{3\pi}\ln(1/\alpha)\approx1040\,\mathrm{MHz}$，与现代 QED 精确值 $1057.833(1)\,\mathrm{MHz}$ 高度一致。兰姆移位的测量是 QED 发展的决定性实验之一。

接下来考察碱金属精细结构。

碱金属的精细结构来自量子亏损对 $j$ 的依赖。钠 D 线分裂：
\[
 \Delta E_{\rm FS}(\mathrm{Na}\,3p)=E(3p_{3/2})-E(3p_{1/2})=17.2\,\mathrm{cm^{-1}}\times hc=2.13\times10^{-3}\,\mathrm{eV},
\]
对应 $\Delta\lambda=0.6\,\mathrm{nm}$（$D_1=589.6\,\mathrm{nm}$, $D_2=589.0\,\mathrm{nm}$）。此分裂远大于氢原子同 $n$ 的精细结构，因为碱金属的有效 $Z_{\rm eff}$ 在核区更大。铯 $6p$ 分裂 $\Delta E=554\,\mathrm{cm^{-1}}$，已接近 $jj$ 耦合区。

接下来考察Thomas 进动的数值估计。

对氢原子 $1s$ 电子，轨道速度 $v\sim\alpha c\sim2.2\times10^6\,\mathrm{m/s}$，轨道频率 $\omega_{\rm orb}\sim4.1\times10^{16}\,\mathrm{rad/s}$。Thomas 进动频率
\[
 \omega_T\sim\frac{v^2}{2c^2}\omega_{\rm orb}\sim\frac{\alpha^2}{2}\omega_{\rm orb}\sim1.1\times10^{12}\,\mathrm{rad/s},
\]
正好是精细结构频率尺度。Thomas 因子 $1/2$ 把「naive 自旋轨道耦合」从 $2\times$ 修正到 $1\times$，是精细结构系数正确的关键。

接下来考察Zitterbewegung 的物理讨论。

自由 Dirac 电子的位置算符含快速振荡项~\eqnrefs{eq:alpha-zitter}，频率 $\omega_Z=2mc^2/\hbar\approx1.55\times10^{21}\,\mathrm{rad/s}$，振幅 $\lambda_C=\hbar/(mc)\approx3.86\times10^{-13}\,\mathrm{m}$（电子 Compton 波长）。这一「颤动运动」来自正负能态的干涉，在 FW 表象中被精确消去——低能电子不表现 Zitterbewegung。但在石墨烯（Dirac 锥有效理论）中，$\omega_Z$ 被压低到可实验观测的尺度（$\sim10^{13}\,\mathrm{rad/s}$），已在输运实验中间接探测。

\begin{derivation}[Bethe 对兰姆移位的对数重整化与 QED 精确计算]
Bethe 1947 年在兰姆移位实验公布后数小时内给出关键估算，其核心是质量重整化与对数截断。分四步重建。

\emph{第一步：电子自能的谱表示。}氢原子中电子的自能修正
\begin{equation*}
  \Delta E_{\rm self}=\langle n|\delta H|n\rangle,\qquad\delta H=\frac{e^2}{2}\int\frac{\dd^3k}{(2\pi)^3}\frac{1}{\omega_k}\left[\frac{1}{H-\omega_k-E_n}-\frac{1}{H_0-\omega_k-mc^2}\right],
\end{equation*}
第一项是束缚态电子与真空光子耦合，第二项是自由电子自能（减去后实现质量重整化）。\(\omega_k=|\vb k|c\) 是光子能量。

\emph{第二步：非相对论近似与对数积分。}对低能态（\(E_n\ll mc^2\)），中间态取非相对论 \(H_0=\vb p^2/(2m)+V\)。Bethe 的关键观察：高能中间态（\(E_{\rm mid}\gtrsim mc^2\)）的贡献被重整化减项消去，剩余积分在 \(\omega_k\in[0,mc^2]\) 上对数发散。引入截断 \(\Lambda=mc^2\)（自然截断来自相对论能区），
\begin{equation*}
  \Delta E_{\rm Lamb}(n,l)\approx\frac{4\alpha^3}{3\pi}\frac{(mc^2)^2}{\hbar^2}\ln\frac{1}{\alpha}\cdot|\psi_{nl}(0)|^2\cdot\delta_{l0}.
\end{equation*}
只有 \(s\) 态（\(\psi_{ns}(0)\ne0\)）有非零移位，\(p\) 态（\(\psi_{np}(0)=0\)）移位为零——与 \(2s_{1/2}-2p_{1/2}\) 分裂一致。对氢 \(2s\)：\(|\psi_{20}(0)|^2=1/(8\pi a_0^3)\)，给出
\begin{equation*}
  \Delta E_{\rm Lamb}(2s)\approx\frac{\alpha^5 mc^2}{6\pi}\ln\frac{1}{\alpha}\approx1040\,\mathrm{MHz}\times h,
\end{equation*}
与实验 \(1057.8\,\mathrm{MHz}\) 在 \(\sim2\%\) 内一致。

\emph{第三步：QED 精确计算。}Bethe 估算的对数精度不足以解释全部移位。完整 QED 处理包括：(i) 精确自能积分（含所有中间态和相对论修正）；(ii) 真空极化图（Uehling 势）；(iii) 反常磁矩修正。精确结果（Bethe--Brown--Stehn 1950，Layzer 1960，Erickson 1977，Mohr 1978）：
\begin{equation*}
  \Delta E_{\rm Lamb}(n,l,j)=\frac{\alpha^3}{\pi n^3}F(n,l,j)\,Z^4\alpha^2 mc^2,
\end{equation*}
其中 \(F\) 是 \(\alpha\) 的幂级数：\(F=A_{40}+A_{41}\ln\alpha+A_{50}\alpha+A_{60}\alpha^2+\cdots\)。主项 \(A_{40}\ln\alpha\) 是 Bethe 对数项，\(A_{41}\) 是 Bethe 常数，\(A_{50}\) 含真空极化和反常磁矩。对氢 \(2s_{1/2}-2p_{1/2}\)：\(A_{40}=4/3\)，精确值 \(1057.833(1)\,\mathrm{MHz}\)。

\emph{第四步：高阶 QED 与质子有限尺寸。}\(A_{60}\) 项（两圈修正）由 Karshenboim 精确计算，贡献 \(\sim-12\,\mathrm{MHz}\)。质子有限尺寸修正 \(\Delta E_{\rm proton}=\frac{2\pi\alpha}{3}|\psi_{ns}(0)|^2 r_p^2\)（\(r_p\) 为质子电荷半径），对 \(2s\) 贡献 \(\sim1.7\,\mathrm{MHz}\)。2010 年 Lamb 移位测量（Pohl et al.，μ氢原子）给出 \(r_p=0.842\,\mathrm{fm}\)，与电子散射值 \(0.877\,\mathrm{fm}\) 的"质子半径疑难"至今未完全解决，是 QED 精度极限下新物理的潜在窗口。

\emph{第五步：Bethe 对数的意义。}Bethe 估算的核心是把紫外发散的对数积分 \(\int\dd\omega/\omega\) 截断在 \(\omega=mc^2\)，给出 \(\ln(1/\alpha)\) 因子。这一对数因子是 QED 微扰展开的标志性结构——在每个能级上，对数来自红外和紫外的"尺度分离"，物理上是电子在原子尺度（\(a_0\)）和 Compton 尺度（\(\lambda_C\)）之间的"无中间尺度"区域的对数贡献。

\emph{Lamb 移位的实验历史。}Lamb--Retherford 1947 年用微波共振测量 \(2s_{1/2}-2p_{1/2}\) 分裂为 \(1057.8\,\mathrm{MHz}\)，直接否定了 Dirac 理论的简并预言，是 QED 发展的关键驱动力。后续测量精度不断提高：Lamb 1952（\(1057.77\,\mathrm{MHz}\)），Robiscoe 1964（\(1057.90\,\mathrm{MHz}\)），Andrews 1971（\(1057.893\,\mathrm{MHz}\)），Hagelberg 1993（\(1057.845\,\mathrm{MHz}\)）。现代测量（Hessels 2000s）精度达 kHz 级，与 QED 理论计算（\(1057.833(1)\,\mathrm{MHz}\)）在 \(10^{-6}\) 相对精度内一致。μ氢原子的 Lamb 移位测量（Pohl 2010，Antognini 2013）给出质子半径值与电子散射的偏差（"质子半径疑难"），是当前 QED 精度极限下新物理搜索的活跃方向。











\emph{[Foldy--Wouthuysen 变换的系统性构造与高阶修正]} Foldy--Wouthuysen (FW) 变换把 Dirac Hamiltonian 中奇-偶耦合项系统消除到指定阶数，给出非相对论极限的精确展开。分四步系统构造。

\emph{第一步：奇偶分解。}Dirac Hamiltonian \(H_D=\beta mc^2+V+\mathcal E+\mathcal O\)，其中 \(\mathcal E\) 是偶算符（与 \(\beta\) 对易，如 \(V\)、\(\boldsymbol p^2/2m\)），\(\mathcal O\) 是奇算符（与 \(\beta\) 反对易，如 \(c\boldsymbol\alpha\cdot\boldsymbol p\)）。偶算符不混合正负能态，奇算符混合——FW 变换的目标是消除奇算符到给定阶数。

\emph{第二步：FW 生成元与迭代。}FW 变换 \(U=\ee^{iS}\)（\(S\) 厄米）由生成元 \(S=-\frac{i\beta\mathcal O}{2mc^2}\) 给出。一次变换后 \(H'=\ee^{iS}H_D\ee^{-iS}\) 的奇算符降阶为 \(O(1/m^2c^4)\)。迭代 \(n\) 次使奇算符降到 \(O(1/m^{2n}c^{4n})\)。Baker--Campbell--Hausdorff 公式 \(H'=\ee^{iS}H\ee^{-iS}=H+i[S,H]+\frac{i^2}{2}[S,[S,H]]+\cdots\) 给出每步的显式展开。

\emph{第三步：非相对论展开。}FW 变换到 \(O(1/m^2c^4)\) 给出
\begin{equation*}
  H_{\rm FW}=mc^2+\frac{\boldsymbol p^2}{2m}+V-\frac{\boldsymbol p^4}{8m^3c^2}+\frac{1}{8m^2c^2}[\boldsymbol\sigma\cdot\boldsymbol p,[\boldsymbol\sigma\cdot\boldsymbol p,V]]+\frac{1}{2m^2c^2}\boldsymbol\sigma\cdot(\nabla V\times\boldsymbol p)+O(1/m^3c^4).
\end{equation*}
各项物理：\(mc^2\) 静止能、\(\boldsymbol p^2/2m\) 动能、\(V\) 外场、\(-\boldsymbol p^4/8m^3c^2\) 相论修正（质量-速度）、Darwin 项 \(\frac{\hbar^2}{8m^2c^2}\nabla^2 V\)（Zitterbewegung）、自旋-轨道 \(\frac{1}{2m^2c^2}\boldsymbol\sigma\cdot(\nabla V\times\boldsymbol p)\)（Thomas 半因子自动出现）。

\emph{第四步：Thomas 半因子的自动出现。}FW 变换中自旋-轨道项的系数 \(\frac{1}{2m^2c^2}\) 含因子 \(1/2\)（Thomas 半因子），它从两个来源的组合自动产生：(1) 正-正投影的 \(\frac{1}{2m^2c^2}\) 和 (2) 奇算符两次作用的 \(-\frac{1}{2}\) 补偿。直接从 Dirac 方程投影到正能态只给出 \(\frac{1}{m^2c^2}\)（缺 Thomas 因子），FW 变换的正确性在于它处理了正负能态的耦合到二阶——这是 Foldy--Wouthuysen 1949 年方法的精髓，把 Thomas 1926 年的运动学修正自然纳入代数框架。
\end{derivation}
